\begin{problem}[When does the whole spectrum have a generic point?]
For a commutative ring $A$, prove the equivalence
\[
\Spec A\text{ is irreducible}
\iff
\sqrt{(0)}\text{ is prime}
\iff
\Spec A\text{ has a generic point}.
\]
Identify that generic point.
\end{problem}

\paragraph{Why this problem matters.}
It corrects the common oversimplification that the generic point is always $(0)$.

\paragraph{Useful ideas before starting.}
Use the irreducibility criterion for $V(I)$ and the fact that $V(0)=\Spec A$.

\paragraph{Guided hints.}
If $\sqrt{(0)}$ is prime, take its closure.

\begin{solution}
The spectrum is
\[
\Spec A=V(0)=V(\sqrt{(0)}).
\]
By the prime/radical criterion for irreducibility,
\[
V(0)\text{ is irreducible}
\iff
\sqrt{(0)}\text{ is prime}.
\]
If $\sqrt{(0)}$ is prime, then it is a point of $\Spec A$ and
\[
\overline{\{\sqrt{(0)}\}}
=
V(\sqrt{(0)})
=
\Spec A,
\]
so it is generic.

Conversely, if $\eta$ is a generic point, then
\[
\Spec A=\overline{\{\eta\}}=V(\eta),
\]
which is irreducible because $\eta$ is prime.  Thus all three conditions are equivalent,
and the unique generic point is
\[
\boxed{\sqrt{(0)}}.
\]
\end{solution}

\paragraph{What happened mathematically.}
Nilpotents may shift the generic point away from the zero ideal without changing the
underlying irreducible topology.

\paragraph{Extensions.}
Show that if $A$ is reduced, then $\Spec A$ is irreducible exactly when $A$ is a domain.

\paragraph{Provenance.}
New problem synthesizing the legacy generic-point theorem with the VI/05 irreducibility
criterion.
